% Chapter 1 — Introduction

\chapter{Introduction}

\label{Chapter1}

Fuel cells convert the chemical energy of hydrogen directly into electricity, with
water and heat as the only by-products. Among them, the proton exchange membrane
fuel cell (PEMFC) is the leading candidate for transport and small stationary
applications thanks to its low operating temperature, high power density, and fast
start-up \cite{pukrushpan2003,daud2017pem}. Their efficient and safe operation,
however, requires feedback control of the gas supply and load: insufficient
reactant flow causes starvation and irreversible degradation, while excessive
anode--cathode pressure difference can damage the membrane \cite{zhu2017,kharrat2025}.
Control design, in turn, requires a model that is simple enough for analysis and
accurate enough to represent the actual hardware --- including the differences
between individual cells of a stack \cite{manana2006}.

This internship, carried out at GIPSA-lab (Universit\'e Grenoble Alpes) under the
supervision of Prof.~Emmanuel Witrant, addresses this need for a small educational
PEMFC stack of ten cells. The objectives are:
\begin{enumerate}[label=\arabic*)]
    \item to develop a simplified, control-oriented dynamic model of a single cell
          and of the stack, grounded in the electrochemical literature;
    \item to identify the model parameters \emph{per cell} from experimental data
          acquired on the bench;
    \item to derive the linearized state-space representation and analyze its
          structural controllability and observability, as a foundation for
          controller and observer design;
    \item to characterize experimentally the dynamic response of the stack to
          resistive load steps, providing the timescales needed for control
          bandwidth selection;
    \item to propose a PI-based control architecture for the PEMFC power system
          through a DC/DC boost converter, following \cite{derbeli2017}.
\end{enumerate}

The report is organized to give a global vision of the work.
Chapter~\ref{Chapter2} summarizes the operating principle of PEMFCs, positions the
work in the literature, and explains the practical constraints that shaped the
control strategy. Chapter~\ref{Chapter3} presents the simplified model, the
parameter-identification procedure, the linearized state-space analysis, and the
experimental protocols. Chapter~\ref{Chapter4} discusses the calibration and
step-response results. Chapter~\ref{Chapter5} presents the control of the PEMFC
power system with a DC/DC boost converter and a PI controller.
Chapter~\ref{Chapter6} concludes and outlines perspectives. The
open-circuit-voltage derivation and the full per-cell parameter table are provided
in Appendices~\ref{app:ocv} and~\ref{app:tables}; the main text is self-contained
and readable without them.
